\documentclass[12pt]{article}
\usepackage{amsmath, amssymb}
\usepackage[margin=1in]{geometry}
\setlength{\parskip}{6pt}
\title{QUBO Formulation: Binary Image Denoising}
\date{}
\begin{document}
\maketitle

\subsection*{Binary Image Denoising}

\paragraph{Problem Statement.}
Given a noisy binary image of dimensions $H \times W$, the objective is to recover a clean binary image by selecting a configuration of pixel labels that minimizes a global energy function. This energy function balances fidelity to the observed noisy pixels against spatial smoothness, penalizing abrupt label changes between adjacent pixels. The task is to find the binary image configuration that achieves the minimum energy.

\paragraph{Decision Variables.}
Let $x_{i,j} \in \{0,1\}$ denote the label of the clean pixel at grid position $(i,j)$, where $i \in \{0, \dots, H-1\}$ and $j \in \{0, \dots, W-1\}$. The variable $x_{i,j}$ takes the value $1$ if the pixel is labeled as foreground (or class $1$), and $0$ otherwise. The total number of decision variables is $N = H \times W$.

\paragraph{QUBO Derivation.}
\textbf{Step 1.} The original energy minimization problem is formulated as:
\begin{align}
\min_{x \in \{0,1\}^{HW}} \mathcal{E}(x) = \sum_{i=0}^{H-1} \sum_{j=0}^{W-1} d_{i,j} x_{i,j} + \sum_{(i,j) \sim (k,l)} w_{i,j,k,l} (x_{i,j} + x_{k,l} - 2 x_{i,j} x_{k,l}),
\end{align}
where $d_{i,j} \geq 0$ is the data cost coefficient for pixel $(i,j)$ derived from the noisy observation, and $w_{i,j,k,l} \geq 0$ is the smoothness penalty weight for neighboring pixels $(i,j)$ and $(k,l)$. The notation $(i,j) \sim (k,l)$ denotes adjacency on the grid (e.g., $4$-connected or $8$-connected).

\textbf{Step 2.} The variables are constrained to the binary domain. The algebraic form of this constraint for each pixel is:
\begin{align}
x_{i,j}^2 - x_{i,j} = 0, \quad \forall i \in \{0, \dots, H-1\}, j \in \{0, \dots, W-1\}.
\end{align}

\textbf{Step 3.} In standard QUBO formulations, constraints are typically enforced via quadratic penalty terms $\lambda (x_{i,j}^2 - x_{i,j})^2$. However, since the domain is explicitly restricted to $\{0,1\}$, the idempotent property $x_{i,j}^2 = x_{i,j}$ holds identically. Substituting this into the penalty form yields $\lambda (x_{i,j} - x_{i,j})^2 = 0$. Consequently, the penalty contribution vanishes, and no explicit penalty terms are required.

\textbf{Step 4.} Combining the objective and the (null) penalty terms, the QUBO Hamiltonian $\mathcal{H}(x)$ is identical to the original energy function:
\begin{align}
\mathcal{H}(x) = \sum_{i=0}^{H-1} \sum_{j=0}^{W-1} d_{i,j} x_{i,j} + \sum_{(i,j) \sim (k,l)} w_{i,j,k,l} (x_{i,j} + x_{k,l} - 2 x_{i,j} x_{k,l}).
\end{align}

\textbf{Step 5.} Expanding $\mathcal{H}(x)$ into the standard QUBO form $\mathcal{H}(x) = \sum_{i,j} Q_{i,j} x_i x_j + c$, where $x_i$ denotes the flattened variable index, we read off the diagonal entries, off-diagonal entries, and offset in closed form:
\begin{align}
Q_{(i,j),(i,j)} &= d_{i,j} + \sum_{(k,l) \sim (i,j)} w_{i,j,k,l}, \\
Q_{(i,j),(k,l)} &= -2 w_{i,j,k,l}, \quad \text{for } (i,j) \sim (k,l), \\
Q_{(i,j),(k,l)} &= 0, \quad \text{for } (i,j) \not\sim (k,l), \\
c &= 0.
\end{align}
Here, the diagonal coefficient for each pixel aggregates its data cost and all connected smoothness weights, while each off-diagonal coefficient corresponds to the interaction between adjacent pixels.

\paragraph{Penalty Weight Justification.}
Because the binary constraints are inherently satisfied by the variable domain, the penalty weight $\lambda$ is effectively zero and no penalty terms are introduced. In scenarios where binary constraints are enforced via penalties, $\lambda$ must be chosen sufficiently large to dominate any feasible objective improvement from violating the constraint. Specifically, $\lambda > \max_{i,j} d_{i,j} + \sum_{(k,l)\sim(i,j)} w_{i,j,k,l}$ would guarantee that infeasible configurations incur a strictly higher energy than any feasible configuration. However, given the explicit restriction $x_{i,j} \in \{0,1\}$, the penalty term identically vanishes, and the Hamiltonian consists solely of the objective terms.

\paragraph{Correctness Argument.}
Minimizing $\mathcal{H}(x)$ over $\{0,1\}^{HW}$ recovers the optimal clean image because the Hamiltonian is constructed to exactly represent the denoising energy function. For any feasible configuration $x \in \{0,1\}^{HW}$, the energy equals the sum of data fidelity and smoothness costs. Since the binary constraints are inherently satisfied, no infeasible configurations exist in the search space. The global minimum of $\mathcal{H}(x)$ therefore corresponds directly to the configuration that minimizes the original energy, guaranteeing that the optimal solution to the QUBO is the optimal denoised image.

\paragraph{Illustrative Example.}
Consider a $2 \times 2$ grid ($H=2, W=2$) with variables $x_{0,0}, x_{0,1}, x_{1,0}, x_{1,1}$. We map them to indices $0, 1, 2, 3$ via the formula $x_{i,j} \to x_{2i+j}$. Assuming $4$-connected neighbors and uniform smoothness weights $w_{i,j,k,l} = w$, the general formulas instantiate to the following concrete QUBO parameters:
\begin{align}
Q_{0,0} &= d_{0,0} + w_{0,0,0,1} + w_{0,0,1,0}, \\
Q_{1,1} &= d_{0,1} + w_{0,1,0,0} + w_{0,1,1,1}, \\
Q_{2,2} &= d_{1,0} + w_{1,0,0,0} + w_{1,0,1,1}, \\
Q_{3,3} &= d_{1,1} + w_{1,1,0,1} + w_{1,1,1,0}, \\
Q_{0,1} &= Q_{1,0} = -2 w_{0,0,0,1}, \\
Q_{0,2} &= Q_{2,0} = -2 w_{0,0,1,0}, \\
Q_{1,3} &= Q_{3,1} = -2 w_{0,1,1,1}, \\
Q_{2,3} &= Q_{3,2} = -2 w_{1,0,1,1}, \\
c &= 0.
\end{align}
The resulting QUBO Hamiltonian is:
\begin{align}
\mathcal{H}(x) &= (d_{0,0} + w_{0,0,0,1} + w_{0,0,1,0}) x_0 + (d_{0,1} + w_{0,1,0,0} + w_{0,1,1,1}) x_1 \nonumber \\
&+ (d_{1,0} + w_{1,0,0,0} + w_{1,0,1,1}) x_2 + (d_{1,1} + w_{1,1,0,1} + w_{1,1,1,0}) x_3 \nonumber \\
&- 2 w_{0,0,0,1} x_0 x_1 - 2 w_{0,0,1,0} x_0 x_2 - 2 w_{0,1,1,1} x_1 x_3 - 2 w_{1,0,1,1} x_2 x_3.
\end{align}
Evaluating $\mathcal{H}(x)$ over all $2^4 = 16$ binary configurations yields the optimal clean image $x^*$ and its minimum energy $\mathcal{H}(x^*)$, which are returned as the solution.

\end{document}